% Source: https://artofproblemsolving.com/community/c3h1523104_how_to_find
% Stable identity: aops-topic-1523104-q01
A point is equally far from $A$ and $B$ exactly when it lies on the perpendicular bisector of $AB$. So the required point (the circumcentre) is the intersection of the perpendicular bisectors of two of the sides; it then also lies on the third.

\textbf{Side $AB$.} Its midpoint is $(4,0)$ and it is horizontal, so its perpendicular bisector is $x=4$.

\textbf{Side $AC$.} Its midpoint is $(1,3)$ and its slope is $3$, so its perpendicular bisector has slope $-\dfrac{1}{3}$:
\[
y-3=-\dfrac{1}{3}(x-1).
\]
At $x=4$ this gives $y=3-1=2$. So the point is $\boxed{O(4,2)}$.

\begin{asy}
usepackage("amsmath");
size(5.5cm);
defaultpen(fontsize(10pt));
pair A=(0,0),B=(8,0),C=(2,6),O=(4,2);
draw(circle(O,sqrt(20)),gray(0.6));
draw(A--B--C--cycle,linewidth(1));
draw((4,-0.9)--(4,6.2),blue+dashed);
draw((0,10/3)--(8,2/3),red+dashed);
real s=0.3; draw((4,0)+(s,0)--(4,0)+(s,s)--(4,0)+(0,s),blue);
dot(O);
label("$A(0,0)$",A,SW);label("$B(8,0)$",B,SE);label("$C(2,6)$",C,NW);label("$O$",O,SW);
label("$x=4$",(4,-0.9),S,blue);
label("\mbox{}",(min(currentpicture,user=true)+max(currentpicture,user=true))/2,black); // keeps the text after the figure black when LaTeX typesets the labels inline
\end{asy}

\textbf{Check.} $OA^2=16+4=20$, $OB^2=16+4=20$, $OC^2=4+16=20$. The common distance is $\sqrt{20}=\boxed{2\sqrt5}$.

(Equivalently, solve $OA^2=OB^2$ and $OA^2=OC^2$ for $O=(x,y)$; the squared terms cancel and two linear equations remain.)
